\section{总结与延伸}

\subsection{核心观点回顾}

本期访谈中，Demis Hassabis 围绕 AGI 的定义、时间线、影响和风险，呈现了一幅完整的图景。以下是核心观点的压缩总结：

\begin{table}[H]
\centering
\begin{tabular}{p{0.25\textwidth} p{0.65\textwidth}}
\toprule
\textbf{议题} & \textbf{核心判断} \\
\midrule
AGI 时间线 & 五年内很有可能实现 \\
AGI 影响规模 & 十倍工业革命、十倍速度 \\
当前 AI 状态 & ``参差不齐的智能''，碎片化是根源 \\
四大缺失能力 & 持续学习、记忆架构、层级规划、一致性推理 \\
Scaling Laws & 收益递减但未终结，竞争转向算法创新 \\
LLM 前景 & 不会被商品化，是 AGI 基础组件 \\
前沿公司差距 & 在拉大，不是在缩小 \\
DeepMind 战略 & 三合一整合，前沿/开源/领域专用三层模型 \\
药物发现 & 5--10 年内变革性成果，目标攻克癌症 \\
首要安全风险 & 恶意使用 > 技术对齐 \\
政策首选 & 建立 AI 版 IAEA \\
能源 & AI 既是问题制造者也是问题解决者 \\
感知偏差 & 短期被高估，长期被严重低估 \\
\bottomrule
\end{tabular}
\caption{Hassabis 核心观点一览}
\end{table}

\subsection{关键语录}

本期访谈中最具冲击力的直接引语：

\begin{enumerate}
    \item \textit{``我有时将 AGI 的到来量化为十倍的工业革命，以十倍的速度发生。''} ——关于 AGI 影响的规模
    \item \textit{``以一种方式提问表现惊人，换个问法在最基础的地方翻车。''} ——关于``参差不齐的智能''
    \item \textit{``支撑现代 AI 的 90\% 的突破，都是由 Google Brain、Research 或 DeepMind 完成的。''} ——关于 DeepMind 的历史贡献
    \item \textit{``在未来五年内实现 AGI 的可能性很大。''} ——关于时间线，与 2010 年 Shane Legg 的 20 年预测一脉相承
    \item \textit{``AI 在短期内被过度炒作，在长期被严重低估。''} ——关于感知偏差
\end{enumerate}

\subsection{对读者最实际的结论}

\begin{enumerate}
    \item \textbf{不要低估 AGI 的时间线}——五年不是虚言，Hassabis 从 2010 年就开始押注这个方向，而且到目前为止他一直在正确的轨道上
    \item \textbf{关注制度适应性}——技术突破越快，制度瓶颈就越凸显。监管改革、教育转型、社会保障体系重设计都需要提前启动
    \item \textbf{AI 安全需要国际协作}——这不是单个公司或单个国家能解决的问题，需要类似 IAEA 的全球机制
    \item \textbf{短期务实、长期敬畏}——不要被炒作裹挟，也不要因为短期的``AI 幻灭''而忽视长期的深刻变革
\end{enumerate}

\subsection{拓展阅读}

\begin{itemize}
    \item Demis Hassabis et al., \textit{AlphaFold: Protein Structure Prediction with Deep Learning}, Nature, 2021
    \item Demis Hassabis et al., \textit{Mastering the Game of Go with Deep Neural Networks and Tree Search}, Nature, 2016
    \item Shane Legg, \textit{Machine Super Intelligence}, PhD Thesis, 2008
    \item Google DeepMind, \textit{Gemini Technical Report}, 2024
    \item Rich Sutton, \textit{The Bitter Lesson}, 2019
    \item Dario Amodei, \textit{Machines of Loving Grace}, 2024
\end{itemize}

\end{document}
